\chapter{Categories, Functors, and Natural Transformations}

\section{Axioms for Categories}
We start by describing categories using \emph{metagraphs}, which do not employ the use of set theory. A metagraph consists of objects, arrows, and two operatons:
\begin{itemize}
	\item Domain, which maps arrows to an object $\operatorname{dom} f=a$;
	\item Codomain, which maps arrows to an object $\operatorname{cod} f=b$.
\end{itemize}
We write $f : a \to b$ to indicate $f$ is an arrow with $\operatorname{dom} f=a$ and $\operatorname{cod} f=b$.

A \emph{metacategory} is a metagraph together with two additional operations:
\begin{itemize}
	\item Identity, assigning each item $a$ an arrow $1_a$;
	\item Composite, assigning each pair $(f,g)$ of arrows with $\operatorname{dom} g=\operatorname{cod} f$, an arrow $g\circ f : \operatorname{dom} f \to \operatorname{cod} g$ called their composite;
\end{itemize}
such that composition is associative and the identity arrows $1_a$ are identities with respect to composition. Alternative notations for the identity $1_a$ include $\mathrm{id} _a$, and simply $a$ itself. Examples of metacategories include the usual category $\Set $ of sets and $\Grp $ of groups, topological spaces with continuous maps as arrows, and compact Hausdorff spaces with the same as arrows.

Since identity arrows are uniquely determined for each object, it's possible to dispense entirely of the objects in a metacategory and define an arrows-only metacategory $C$. There exist a collection of arrows, for which some pairs $(f,g)$ of arrows are called \emph{composable}. For any composable arrows, we say the composition $g\circ f$ is defined, and we can then define an identity arrow to be any arrow $u$ satisfying $f\circ u=f$ and $u\circ g=u$ for any $f,g$ such that the composition is defined. The data must then satisfy the axioms
\begin{itemize}
	\item $(k\circ g)\circ f$ is defined iff $k\circ (g\circ f)$ is defined, and they are equal;
	\item The composite $kgf$ is defined if and only if the composites $kg$ and $gf$ are defined;
	\item For each arrow $g$ of $C$, there exist identity arrows $u,u'$ such that $gu$ and $u'g$ are defined.
\end{itemize}

The final axiom indeed shows that identity arrows are unique in this regard, and allows defining of a domain and codomain for the arrow $g$. An arrows-only metacategory thus satisfies the usual metacategory axioms, if identity arrows are taken as objects.

\section{Categories}

\begin{definition}[Directed Graph]\label{dfn:1.1}
	A directed graph is a set $O$ of objects, a set $A$ of arrows, and two functions
	\[\begin{tikzcd}[row sep = 2.5em, column sep = 2.5em]
		A\ar[r, shift left, "\text{dom} "]\ar[r, shift right, "\text{cod} "']&O
	\end{tikzcd}\]
	where the set of composable pairs is the set
	\[
		A\times_OA=\left\{ (g,f)\mid g,f\in A,\operatorname{dom} g=\operatorname{cod} f \right\} 
	.\]
\end{definition}

\begin{definition}[Category]\label{dfn:1.2}
	A category is a directed graph with two additional functions
	\[\begin{tikzcd}[row sep = 2.5em, column sep = 2.5em]
		O\ar[r, "\mathrm{id} "]&A,&A\times _OA\ar[r, "\circ "]&A,\\
		c\ar[r, mapsto]&1_c,&(g,f)\ar[r, mapsto]&g\circ f,
	\end{tikzcd}\]
	where $g\circ f$ is usually written $gf$, such that
	\[
		\operatorname{dom} 1_a=\operatorname{cod} 1_a,\qquad \operatorname{dom} gf=\operatorname{dom} f,\qquad \operatorname{cod} gf=\operatorname{cod} g
	.\]
	for all $a\in O$ and $(g,f)\in A\times _OA$, such that a metacategory is formed.
\end{definition}

That is, a category is an interpretation of a metacategory within set theory. If $C$ is a category, we usually write $a\in C$ or $f\in C$ to say $a$ is an object of $C$, or $f$ is an arrow of $C$. We also write
\[
	\Hom(b, c) =\left\{ f\in C\mid \operatorname{dom} f=b,\operatorname{cod} f=c \right\} 
\]
for the set of arrows $b \to c$. We will show later that this definition of a category is the same as saying a category is a monoid for the product $\times_O$.

Examples of categories include
\begin{itemize}
	\item $\mathbf{0} $, the empty category;
	\item $\mathbf{1} $, the category with one object and one arrow;
	\item $\mathbf{2} $, the category with two objects and precisely one non-identity arrow $a\to b$;
	\item $\mathbf{3} $, the category with three objects whose non-identity arrows form a triangle
		\[\begin{tikzcd}[row sep = 2.5em, column sep = 2.5em]
			&\bullet\ar[dr] \\
			\bullet \ar[ur] &&\bullet \ar[ll] 
		\end{tikzcd}\]
	\item $\downdownarrows$, the category with two objects and two non-identity arrows $a\rightrightarrows b$; we call such arrows \emph{parallel arrows}.
\end{itemize}

A category is called \emph{discrete} when all arrows are identities. Any set forms a discrete category, and any discrete category is completely determined by it's arrows.

A \emph{monoid} is a category with one object. A monoid is thus fully described by it's arrows, and it's composition operation. A monoid can thus be described as a set $M$ (of arrows) together with a binary operation $M\times M\to M$ and an identity element; such an object is precisely a semigroup with unity. Thus, for any category $C$ and any object $a\in C$, the category $\Hom(a, a) $ forms a monoid. In this vein of thought, a group is thus a category where every arrow has a two-sided inverse.

Given a set $V$ of sets, the category $\mathbf{Ens}_V$ takes $V$ as the objects, and morphisms as any regular functions between sets in $V$. $\mathbf{Ens} $ refers to an arbitrary category of this type.

An ordinal number $n$ can be seen as the set $\left\{ 0, \ldots , n-1 \right\} $, and hence each ordinal is a category when viewed as a preorder. We can also consider the category $\omega $, the first infinite cardinal, which when viewed as a category looks like
\[\begin{tikzcd}[row sep = 2.5em, column sep = 2.5em]
	0\ar[r] &1\ar[r] &2\ar[r] &3\ar[r] &\ldots 
\end{tikzcd}\]

The category $\Delta $ takes all finite ordinals as objects and all order-preserving functions $f : m \to n$ as arrows. This category is generally called the \emph{simplical category}, and it is extremely important.

We want an actual category $\Set $, which is clearly not a set itself. We will circumvent this by allowing there to exist a ``universe'' $\mathcal{U} $, and will call a set ``small'' if it belongs to the universe. We will construct this explicitly later (and I will be calling small sets ``sets'' and large sets ``classes'' like a normal fucking person). This allows us to construct
\begin{itemize}
	\item $\Set $, objects are sets and morphisms are functions;
	\item $\Set_*$, objects are sets with a selected point, and morphisms are functions preserving these points;
	\item $\mathbf{Ens} $, category of all classes and functions within a variable set $V$;
	\item $\mathbf{Cat} $, category of small categories and morphisms as functors;
	\item $\Grp $, category of groups and morphisms as group homomorphisms;
	\item $\mathbf{Ab} $, category of abelian groups;
	\item $\mathbf{Mon} $, category of (small) monoids;
	\item $\Ring $, category of rings;
	\item $\mathbf{CRing} $, category of commutative rings;
	\item $R$-$\mathbf{Mod} $, category of left $R$-modules;
	\item $\mathbf{Mod}$-$R$, category of right $R$-modules;
	\item $\mathbf{Top} $, category of topological spaces;
	\item $\mathbf{Toph} $, category of topological spaces with morphisms as homotopy classes of maps;
	\item $\mathbf{Top}_*$, category of topological spaces with a distinguised base point; base preserving continuous maps.
\end{itemize}

\section{Functors}

\begin{definition}[Functor]\label{dfn:1.3}
	Let $C,D$ be categories. A functor $T : C \to D$ consists of an object function $T$ mapping each $c\in C$ to some $Tc\in D$, and an arrow function mapping each $f : c \to c'$ in $C$ to $Tf : Tc \to Tc'$ in $D$, in such a way that
	\[
		T1_c=1_{Tc} ,\qquad T(g\circ f)=T(g)\circ T(f)
	\]
	whenever the composite $g\circ f$ is defined in $C$.
\end{definition}

For example, the power set functor $\mathscr{P} : \Set  \to \Set $, defined such that for any $f : A \to B$, the new map $\mathscr{P} f : \mathscr{P} A \to \mathscr{P} B$ is such that for any $S\subseteq A$, $\mathscr{P} f(S)=f(S)$.

Functors first arose in algebraic topology. For any $n$, a topological space $X$ can be assigned an abelian group $H_n(X)$ called the $n$-th homology group of $X$. However, for any $: X \to Y$, there exists an induced $H_n(f): H_n(X) \to H_n(Y)$ morphism of groups, and thus $H_n$ becomes a functor $\Top \to \mathbf{Ab} $. We can in fact show homologygroups are invariant under homotopy, and so we can equally define a functor $\mathbf{Toph} \to \Grp $. The homotopy groups $\pi _n$ are also functors; since they depend on a base point they define a functor $\Top_* \to \Grp $.

A functor which ``forgets'' underlying structure of an object is called a \emph{forgetful functor}. For example, there is an apparent forgetful functor $U : \Grp  \to \Set $, another $U : \Ring  \to \mathbf{Ab}  $, and so on.

Clearly functors may be composed, and moreover there is an identity functor $1_C$ for any category $C$. We may thus consider the metacategory of all categories, or the category $\mathbf{Cat} $ of all small categories.

A functor is called an \emph{isomorphism} if it is two-sided invertible; in general isomorphism is too strong a condition, and weaker conditions are more useful. A functor is \emph{full} if for every $c,c'\in C$ and every arrow $g : Tc \to Tc'$ in $D$, there exists $f : c \to c'$ in $C$ such that $g=Tf$. A functor is \emph{faithful}, or an embedding, if for all $c_1,c_2\in C$ and parallel arrows $f_1,f_2: c_1 \to c_2$ such that $Tf_1=Tf_2$, we have $f_1=f_2$.

These properties can be realized more intuitively. A functor defines functions $T : \Hom(c, c')  \to \Hom(Tc, Tc') $. A functor is full when these functions are surjective; it is faithful when they are injective. A fully faithful functor is then a functor where these functions are bijections.

A subcategory $S$ of a category $C$ is a subset of the objects and arrows in $C$ in such a way that $S$ forms a category. There is automatically a faithful inclusion functor $S\hookrightarrow C$; we say a subcategory is \emph{full} when the inclusion functor is full. Intuitively, $S$ contains all arrows in $C$ for the specified subcollection of objects.

\section{Natural Transformations}
\begin{definition}[Natural Transformation]\label{dfn:1.4}
	Let $T,S : C \to D$ be functors. A natural transformation $\tau : S \Rightarrow T$ is a function which assigns to each $c\in C$ an arrow $\tau _c=\tau c : Sc \to Tc$ of $D$, called the \emph{component of $\tau $ at $c$}, such that for all $f : c \to c'$ in $C$, the diagram
	\[\begin{tikzcd}[row sep = 2.5em, column sep = 2.5em]
		Sc\ar[" \tau c ", r] \ar[" Sf "', d] &Tc\ar[" Tf ", d] \\
		Sc'\ar[" \tau c' ", r] &Tc'
	\end{tikzcd}\]
	commutes.
\end{definition}

We say any component $\tau c : Sc \to Tc$ is \emph{natural} in $c$. We can think of a natural transformation $\tau : S \Rightarrow T$ as a way of transforming the image of a functor $S$ into that of a functor $T$. For example, given a triangle
\[\begin{tikzcd}[row sep = 2.5em, column sep = 2.5em]
	a \ar[" f ", dr] \ar[" h "', dd] \\
	&b \ar[" g ", dl] \\
	c
\end{tikzcd}\]
we would expect the diagram
\[\begin{tikzcd}[row sep = 2.5em, column sep = 2.5em]
	Sa\ar[" \tau a ", rrr] \ar[" Sf ", dr] \ar[" Sh ", dd] &&&Ta\ar[" Tf ", dr] \ar[" Th ", dd] \\
							       &Sb \ar[" Sg ", dl] &&&Tb \ar[" Tg ", dl] \\
	Sc\ar[" \tau c "', rrr] &&&Tc
	\arrow[from=2-2, to=2-5, crossing over]
\end{tikzcd}\]
commute. Indeed, it does.

Natural transformations define \emph{morphisms of functors}. If a natural transformation has every component invertible, we call it a \emph{natural isomorphism} and write $S\cong T$.

For example, the determinant is a natural transformation $\det_K : \operatorname{GL}_n \Rightarrow \left(  \right) ^*$, where $K\mapsto K^*$ maps a ring to it's group of units. These are clearly two functors 
